% !TeX root = ../../Book1.tex
\section{定义}
\label{b1:sec:2101}

上一章把导数关系写成测试积分恒等式。现在给函数及其弱导数同时加上可积性要求，就能用一个范数度量它们。这种做法适合处理取极限的问题：函数可以不再光滑，但只要各阶弱导数仍有适当的积分控制，极限就有希望留在同一个空间中。

\ProseParagraphBreak

以下固定非空开集 \(\Omega\subset\mathbb R^d\)，使用 Lebesgue 测度，标量域为 \(\mathbb K=\mathbb R\) 或 \(\mathbb C\)。不预设 \(\Omega\) 有界、连通或边界光滑。弱导数始终按\cref{b1:def:weak-derivative}理解；函数及其导数均按几乎处处相等取等价类。
% 实读来源：Hunter2001Applied作者公开第12章印刷362--364页，定义12.66及其后的范数、内积。
% 此处完整规定本书多指标范数；对局部空间、边界可积性及高阶计数另作说明。

\subsection{\texorpdfstring{\(W^{1,p}\)}{W1,p}}
\label{b1:subsec:210101}

\begin{definition}\label{b1:def:sobolev-first-order}
给定 \(1\le p\le\infty\)，若 \(u\in L^p(\Omega)\)，且它的每个一阶弱偏导数 \(\partial_i u\) 均存在并属于 \(L^p(\Omega)\)，则称 \(u\) 属于一阶\term[sobolev-kongjian]{Sobolev 空间}[Sobolev space]，记为 \(u\in W^{1,p}(\Omega)\)。规定
\[
 \|u\|_{W^{1,p}(\Omega)}
 \coloneqq\left(\|u\|_{L^p(\Omega)}^p+
         \sum_{i=1}^d\|\partial_i u\|_{L^p(\Omega)}^p\right)^{1/p}
 \quad(1\le p<\infty),
\]
\[
 \|u\|_{W^{1,\infty}(\Omega)}
 \coloneqq\max\bigl\{\|u\|_{L^\infty(\Omega)},
                \|\partial_1u\|_{L^\infty(\Omega)},\ldots,
                \|\partial_du\|_{L^\infty(\Omega)}\bigr\}.
\]
\end{definition}
\symidx{\(W^{1,p}(\Omega)\)：函数及一阶弱导数属于\(L^p\)的空间}

这些弱导数首先确实可以按上一章定义。因为对每个紧集 \(K\Subset\Omega\)，Hölder 不等式给出
\[
 \int_K|u|\,\mathrm dx
 \le m_d(K)^{1-1/p}\|u\|_{L^p(\Omega)}
 \quad(1\le p<\infty),
\]
而 \(p=\infty\) 时右端换为 \(m_d(K)\|u\|_\infty\)。因此 \(L^p(\Omega)\subset L^1_{\mathrm{loc}}(\Omega)\)，即使 \(\Omega\) 的总体积无限也成立。\secref{b1:sec:2004}的唯一性又保证范数不依赖弱导数代表的选择。

\ProseParagraphBreak

定义中既控制导数，也控制函数本身。在有界开集上，非零常函数的梯度为零，函数本身却不是零等价类；因此只用 \(\|\nabla u\|_p\) 不能定义这里的范数。涉及边界条件或把常数取商的空间，需要以后另作定义。

\ProseParagraphBreak

在 \((-1,1)\) 上，\(u(x)=|x|\) 属于每个 \(W^{1,p}\)，包括 \(p=\infty\)，尽管它在原点不可微。相反，一个函数即使在开集内处处光滑，也不保证满足全域积分条件。例如 \(u(x)=\sqrt{x}\) 在 \((0,1)\) 内光滑，且函数本身有界，但
\[
 \int_0^1|u'(x)|^p\,\mathrm dx
 =2^{-p}\int_0^1x^{-p/2}\,\mathrm dx
\]
只在 \(p<2\) 时有限。因此它属于 \(W^{1,p}(0,1)\) 当且仅当 \(1\le p<2\)。这里的障碍发生在趋近边界时，不在任何内部点。

\subsection{\texorpdfstring{\(W^{k,p}\)}{Wk,p}}
\label{b1:subsec:210102}

令 \(k\in\mathbb N_0\)，记
\[
 A_k\coloneqq\{\alpha\in\mathbb N_0^d:|\alpha|\le k\},
 \qquad N_k\coloneqq\#A_k=\binom{d+k}{k}.
\]
约定 \(\partial^0u=u\)。对高阶空间，需要每一个不超过 \(k\) 阶的弱导数，而不只是某几个最高阶组合。

\begin{definition}\label{b1:def:sobolev-higher-order}
若 \(u\in L^p(\Omega)\)，且对每个 \(\alpha\in A_k\)，弱导数 \(\partial^\alpha u\) 都属于 \(L^p(\Omega)\)，则称 \(u\in W^{k,p}(\Omega)\)。其范数规定为
\begin{equation}\label{b1:eq:sobolev-norm}
 \|u\|_{W^{k,p}(\Omega)}
 \coloneqq\left(\sum_{\alpha\in A_k}
              \|\partial^\alpha u\|_{L^p(\Omega)}^p\right)^{1/p}
 \quad(p<\infty),
\end{equation}
而 \(p=\infty\) 时取这些导数的 \(L^\infty\) 范数的最大值。特别地，\(W^{0,p}(\Omega)=L^p(\Omega)\)。
\end{definition}
\symidx{\(W^{k,p}(\Omega)\)：函数及所有不超过\(k\)阶的弱导数均属于\(L^p\)的空间}

对整数 \(m\geq1\)，正式约定
\[
 W^{k,p}(\Omega;\mathbb R^m)
 \coloneqq\bigl\{u=(u^1,\ldots,u^m):u^a\in W^{k,p}(\Omega),\ 1\leq a\leq m\bigr\}.
\]
这里
\[
 \partial^\alpha u
 \coloneqq\bigl(\partial^\alpha u^1,\ldots,\partial^\alpha u^m\bigr),
\]
并固定用 \(\mathbb R^m\) 的欧氏范数计算每个 \(\|\partial^\alpha u\|_{L^p(\Omega;\mathbb R^m)}\)。向量值 Sobolev 范数仍按\eqref{b1:eq:sobolev-norm}对 \(\alpha\in A_k\) 作 \(p\) 次和；当 \(p=\infty\) 时，取这些向量值 \(L^\infty\) 范数的最大值。这一约定固定了后文定量估计中的有限维范数，而不只依赖有限维范数等价性。
\symidx{\(W^{k,p}(\Omega;\mathbb R^m)\)：各分量属于\(W^{k,p}(\Omega)\)的有限维向量值 Sobolev 空间}

\ProseParagraphBreak

若 \(|\alpha|+|\beta|\le k\)，则 \(\partial^\alpha(\partial^\beta u)=\partial^{\alpha+\beta}u\)。这不是另加一条经典可微假设：等式先由分布导数可交换得到，再由局部可积代表的唯一性成为弱导数的等式。因此也可以递归地检查 \(u\) 及其低阶导数是否具有下一阶弱导数。

\ProseParagraphBreak

同一空间常配有不同但等价的范数。例如令
\[
 S(u)=\sum_{\alpha\in A_k}\|\partial^\alpha u\|_p.
\]
对 \(p<\infty\)，有限维 Hölder 不等式给出
\[
 \|u\|_{W^{k,p}}\le S(u)
 \le N_k^{1-1/p}\|u\|_{W^{k,p}},
\]
对 \(p=\infty\) 则有 \(\|u\|_{W^{k,\infty}}\le S(u)\le N_k\|u\|_{W^{k,\infty}}\)。所以空间及收敛概念相同，但后续若要追踪精确常数，须回到本书的\eqref{b1:eq:sobolev-norm}，不能把等价范数当作数值相等。

\ProseParagraphBreak

当 \(k\ge1\) 时，记最高阶部分为
\[
 |u|_{W^{k,p}(\Omega)}
 \coloneqq\left(\sum_{|\alpha|=k}\|\partial^\alpha u\|_p^p\right)^{1/p}
 \quad(p<\infty),
\]
无穷端点仍取最大值。这一般只是半范数。比如在有界域上，次数低于 \(k\) 的多项式具有零的最高阶半范数，却未必为零函数。

\ProseParagraphBreak

提高 \(k\) 会缩小空间：\(W^{k+1,p}(\Omega)\subset W^{k,p}(\Omega)\)，且本书范数下包含映射的范数不超过 \(1\)。提高 \(p\) 则还要看域的测度。若 \(m_d(\Omega)<\infty\) 且 \(1\le p\le q\le\infty\)，逐项使用 Hölder 可得
\[
 \|u\|_{W^{k,p}(\Omega)}
 \le N_k^{1/p-1/q}\cdot m_d(\Omega)^{1/p-1/q}
       \|u\|_{W^{k,q}(\Omega)}.
\]
在无限测度域上没有这条无条件包含关系；常函数 \(1\) 在 \(\mathbb R^d\) 上属于 \(W^{k,\infty}\)，却不属于任何有限 \(p\) 的 \(W^{k,p}\)。

\subsection{\texorpdfstring{\(H^k=W^{k,2}\)}{Hk=Wk,2}}
\label{b1:subsec:210103}

\begin{definition}\label{b1:def:sobolev-hilbert}
将 \(W^{k,2}(\Omega)\) 记作 \(H^k(\Omega)\)，并定义
\begin{equation}\label{b1:eq:sobolev-inner-product}
 \langle u,v\rangle_{H^k(\Omega)}
 \coloneqq\sum_{\alpha\in A_k}\int_\Omega
       \partial^\alpha u\,\overline{\partial^\alpha v}\,\mathrm dx.
\end{equation}
实值情形省略复共轭。
\end{definition}
\symidx{\(H^k(\Omega)\)：\(W^{k,2}(\Omega)\)及其导数内积}

内积对第一变量线性，与\chapref{b1:ch:08}约定一致。零阶项保证正定性，并且 \(\langle u,u\rangle_{H^k}=\|u\|_{W^{k,2}}^2\)。因此 \(H^k\) 不仅表示另一种记号，还突出了二次能量与正交方法可以使用的结构；要称它为 Hilbert 空间，仍须在下一节证明完备性。

\ProseParagraphBreak

这里每个多重指标只计一次。例如二维 \(H^2\) 范数含有一个 \(\|\partial_1\partial_2u\|_2^2\) 项；若把 Hessian 矩阵每个位置都逐项平方求和，混合导数会计两次，得到的是另一个等价范数。两种约定都常用，本文需要数值恒等式时始终以\eqref{b1:eq:sobolev-inner-product}为准。

\ProseParagraphBreak

Hunter 与 Nachtergaele 的《Applied Analysis》\cite[Section 12.9, Definition 12.66 and the following discussion]{Hunter2001Applied}把这些空间放在弱导数之后引入。接下来我们将把导数作为 \(L^p\) 中的有限组分量，直接从\chapref{b1:ch:10}的完备性推导其基本结构。

\subsection{局部 Sobolev 空间}
\label{b1:subsec:210104}

全域可积性会同时检测函数在边界附近及无穷远处的行为。研究开集内部的运算时，往往只需要把这些位置暂时排除。

\begin{definition}\label{b1:def:sobolev-local}
若 \(u\in L^1_{\mathrm{loc}}(\Omega)\)，且每个 \(|\alpha|\le k\) 阶弱导数都存在并属于 \(L^p_{\mathrm{loc}}(\Omega)\)，则称 \(u\) 属于\term[jubu-sobolev-kongjian]{局部 Sobolev 空间}[local Sobolev space]，记作 \(W^{k,p}_{\mathrm{loc}}(\Omega)\)。这里 \(L^p_{\mathrm{loc}}\) 表示在每个紧子集上 \(p\) 次幂可积，或在 \(p=\infty\) 时本性有界。
\end{definition}
\symidx{\(W^{k,p}_{\mathrm{loc}}(\Omega)\)：局部Sobolev空间}

向量值局部空间也逐分量定义：
\[
 W^{k,p}_{\mathrm{loc}}(\Omega;\mathbb R^m)
 \coloneqq\bigl\{u=(u^1,\ldots,u^m):u^a\in W^{k,p}_{\mathrm{loc}}(\Omega),\ 1\leq a\leq m\bigr\}.
\]
因此后文的向量值链式法则不再把“逐分量理解”作为未写出的定义。

\ProseParagraphBreak

等价地，对每个开集 \(V\) 满足 \(\overline V\Subset\Omega\)，都有 \(u|_V\in W^{k,p}(V)\)。一个方向由限制立即得到。反过来，用可数个这类 \(V\) 覆盖 \(\Omega\)；重叠部分的弱导数由唯一性相等，故能拼接为同一个局部 \(L^p\) 函数，再用\cref{b1:prop:weak-derivative-locality}得到整个 \(\Omega\) 上的弱导数关系。

\ProseParagraphBreak

例如 \(x\mapsto1/x\) 属于 \((0,1)\) 上的每个 \(W^{k,p}_{\mathrm{loc}}\)，但并不属于该区间上的 \(L^1\)。另一个方向，\(W^{k,p}(\Omega)\) 总包含于 \(W^{k,p}_{\mathrm{loc}}(\Omega)\)，而局部成员只有在所有导数的全域范数也有限时才回到全域空间。

\ProseParagraphBreak

局部收敛同样按内部区域理解：\(u_j\to u\) 于 \(W^{k,p}_{\mathrm{loc}}(\Omega)\)，是指对上述每个 \(V\)，都有 \(\|u_j-u\|_{W^{k,p}(V)}\to0\)。只要选一列内部开集穷竭，使每个紧子集最终落在其中，检查这可数列就足够。局部化并不规定边界值，更不意味着把函数在域外补零后仍属于同阶 Sobolev 空间；上一章的端点点质量已经说明零延拓可能产生额外的分布导数。
